% ============================================================
%  PAUTA PEP 1 — FISICA I (FORMA B) — SOLO PROBLEMA 1
%  Usach Premium
%  Compila con: pdflatex pauta_pep1_forma_b_p1.tex  (2 veces)
% ============================================================

\documentclass[letterpaper, 11pt]{article}

\usepackage[utf8]{inputenc}
\usepackage[spanish]{babel}
\usepackage{amsmath, amssymb}
\usepackage{geometry}
\usepackage{fancyhdr}
\usepackage{enumitem}
\usepackage{graphicx}
\usepackage{hyperref}
\usepackage{parskip}
\usepackage{xcolor}
\usepackage{tcolorbox}
\usepackage{eso-pic}
\usepackage{tikz}
\usetikzlibrary{patterns, arrows.meta, decorations.pathmorphing, calc, babel}
\tcbuselibrary{skins, breakable}

\geometry{letterpaper, margin=2cm, top=3cm, bottom=2.5cm, headheight=2cm}

\definecolor{celesteSuave}{HTML}{F0F7FF}
\definecolor{azulFuerte}{HTML}{1A5276}
\definecolor{verdePaso}{HTML}{1E8449}
\definecolor{rojoAlerta}{HTML}{C0392B}
\definecolor{naranjaAcento}{HTML}{D35400}
\definecolor{enlace}{HTML}{1A5276}

\newcommand{\logoup}{./images/logo_up.png}
\newcommand{\logousach}{./images/logo_usach.png}
\newcommand{\logofondo}{./images/logo_up.png}

\newcommand{\institucion}{Usach Premium}
\newcommand{\curso}{Fisica I para Ingenieria}
\newcommand{\guiasnumero}{Pauta PEP 1 --- Forma B --- Problema 1}
\newcommand{\semestre}{1er. Semestre de 2026}
\newcommand{\preparadopor}{Usach Premium.}
\newcommand{\webinstitucion}{www.usachpremium.cl}
\newcommand{\instagraminstitucion}{https://www.instagram.com/usach.premium}
\newcommand{\transparencia}{0.04}
\newcommand{\fraseportada}{"Por y para estudiantes"}

\hypersetup{colorlinks=true, linkcolor=enlace, urlcolor=enlace,
            citecolor=enlace, pdfborder={0 0 0}}

\pagestyle{fancy}
\fancyhf{}
\fancyhead[L]{\includegraphics[height=1.2cm]{\logoup}}
\fancyhead[R]{%
    \begin{minipage}[b]{0.45\textwidth}
        \raggedleft\fontsize{9pt}{11pt}\selectfont
        \textbf{\color{azulFuerte}\institucion} \\
        \curso\ --- \guiasnumero \\
        \textit{\color{gray}@usach.premium}
    \end{minipage}\hspace{0.1cm}%
    \includegraphics[height=1.2cm]{\logousach}}
\fancyfoot[L]{\fontsize{9pt}{11pt}\selectfont \textit{\fraseportada}}
\fancyfoot[R]{\fontsize{9pt}{11pt}\selectfont P\'agina \thepage}
\renewcommand{\headrulewidth}{0.4pt}
\renewcommand{\footrulewidth}{0.4pt}
\fancypagestyle{plain}{\fancyhf{}\renewcommand{\headrulewidth}{0pt}
                                   \renewcommand{\footrulewidth}{0pt}}

\newcommand{\MarcaAgua}{
    \AddToShipoutPicture{
        \AtPageCenter{
            \begin{tikzpicture}[remember picture, overlay]
                \node[opacity=\transparencia] at (0,0)
                    {\includegraphics[width=0.7\textwidth]{\logofondo}};
            \end{tikzpicture}}}}

\newtcolorbox{paso}[2][]{
    colback=celesteSuave, colframe=azulFuerte, fonttitle=\bfseries,
    coltitle=white, colbacktitle=azulFuerte,
    title={\large Paso #2},
    arc=6pt, boxrule=1pt, enhanced, breakable, #1
}
\newtcolorbox{justif}[1]{
    colback=naranjaAcento!8, colframe=naranjaAcento, fonttitle=\bfseries\small,
    title={#1}, arc=5pt, boxrule=0.8pt, enhanced, breakable
}
\newtcolorbox{enunciadobox}[1]{
    colback=azulFuerte!10, colframe=azulFuerte, fonttitle=\bfseries,
    coltitle=white, colbacktitle=azulFuerte, title={#1},
    arc=7pt, boxrule=1.2pt, enhanced, breakable
}
\newtcolorbox{resultado}{
    colback=azulFuerte, colframe=azulFuerte, colupper=white,
    arc=10pt, boxrule=1.5pt, enhanced, drop shadow={black!15}
}

\begin{document}
\thispagestyle{plain}
\MarcaAgua

\AddToShipoutPicture*{%
    \put(54,700){\includegraphics[width=2cm]{\logoup}}
    \put(420,706){\includegraphics[width=5cm]{\logousach}}}

\vspace*{0.5cm}
\begin{center}
    \color{azulFuerte}
    {\huge \textbf{PAUTA DE CORRECCI\'ON}} \\[0.4cm]
    {\LARGE \textbf{PEP 1 --- \curso}} \\[0.3cm]
    {\Large Forma B \quad$\bullet$\quad Problema 1} \\[0.6cm]
    {\footnotesize\color{gray} Considera $g=9{,}8\ \mathrm{m/s^2}$. Resultados finales aproximados a 2 decimales.}
\end{center}

\vspace{0.4cm}

\begin{enunciadobox}{Problema 1 --- Cami\'on de reparto aut\'onomo \hfill (20 puntos)}
Cami\'on de masa $M=1052\,\mathrm{kg}$ que se mueve seg\'un el gr\'afico $v$-$t$
(tres etapas). Un vehículo de apoyo ($m=2510\,\mathrm{kg}$) parte del mismo
origen con retraso de $2\,\mathrm{s}$, rapidez inicial $v_0$ desconocida y
aceleraci\'on $\vec{a}=(3{,}0\,\mathrm{m/s^2})\,\hat{\imath}$. Ambos se encuentran
en la \textbf{segunda etapa} y en ese instante el apoyo alcanza $v_f=26\,\mathrm{m/s}$.
\end{enunciadobox}

\medskip
\begin{center}
\begin{tikzpicture}[x=0.28cm,y=0.13cm]
  \draw[-{Latex[length=2mm]},line width=0.6pt] (0,0) -- (30,0) node[below left=1pt]{$t\,(\mathrm{s})$};
  \draw[-{Latex[length=2mm]},line width=0.6pt] (0,0) -- (0,32) node[right=2pt]{$v\,(\mathrm{m/s})$};
  \draw[dashed] (0,24) -- (26,24);
  \draw[dashed] (8,0) -- (8,24);
  \draw[dashed] (20,0) -- (20,24);
  \draw[line width=1pt, azulFuerte] (0,0) -- (8,24) -- (20,24) -- (26,0);
  \node[below] at (8,0) {$8$}; \node[below] at (20,0) {$20$}; \node[below] at (26,0) {$26$};
  \node[left] at (0,24) {$24$}; \node[below left] at (0,0){$0$};
  \node[verdePaso] at (4,13) {\scriptsize Etapa 1};
  \node[verdePaso] at (14,26) {\scriptsize Etapa 2};
  \node[verdePaso] at (24,13) {\scriptsize Etapa 3};
\end{tikzpicture}
\end{center}

% ------------------- P1 (a) -------------------
\begin{paso}{1.a: Aceleraciones $a_1$ y $a_3$ (pendiente del gr\'afico)}
La aceleraci\'on es la pendiente de la recta $v$-$t$ en cada tramo.

\textbf{Etapa 1} ($0\le t\le 8$): la velocidad pasa de $0$ a $24\,\mathrm{m/s}$:
\[
a_1=\frac{\Delta v}{\Delta t}=\frac{24-0}{8-0}=3\ \mathrm{m/s^2}.
\]
\textbf{Etapa 3} ($20\le t\le 26$): la velocidad pasa de $24$ a $0\,\mathrm{m/s}$:
\[
a_3=\frac{0-24}{26-20}=\frac{-24}{6}=-4\ \mathrm{m/s^2}.
\]
\begin{resultado}\centering
$\displaystyle a_1 = 3{,}00\ \mathrm{m/s^2}\qquad\text{y}\qquad a_3 = -4{,}00\ \mathrm{m/s^2}$
\end{resultado}
\end{paso}

% ------------------- P1 (b) -------------------
\begin{paso}{1.b: Ecuaciones de posici\'on $x(t)$ (origen en la posici\'on inicial del cami\'on)}
Se usa $x(t)=x_0+v_0 t+\tfrac12 a t^2$ en cada tramo, encadenando posici\'on y
velocidad al final de cada etapa (la etapa 2 es MRU con $v=24\,\mathrm{m/s}$).

\textbf{Etapa 1} ($0\le t\le 8$): parte del reposo en el origen, $a_1=3$:
\[
x_1(t)=\tfrac12(3)t^2=1{,}5\,t^2,\qquad x_1(8)=96\ \mathrm{m}.
\]
\textbf{Etapa 2} ($8\le t\le 20$): parte de $96\,\mathrm{m}$ con $v=24$ constante:
\[
x_2(t)=96+24\,(t-8),\qquad x_2(20)=384\ \mathrm{m}.
\]
\textbf{Etapa 3} ($20\le t\le 26$): parte de $384\,\mathrm{m}$ con $v=24$ y $a_3=-4$:
\[
x_3(t)=384+24\,(t-20)-2\,(t-20)^2,\qquad x_3(26)=456\ \mathrm{m}.
\]
\begin{resultado}\centering
$\displaystyle x_1(t)=1{,}5\,t^2\quad;\quad x_2(t)=96+24(t-8)\quad;\quad x_3(t)=384+24(t-20)-2(t-20)^2$
\end{resultado}
\end{paso}

% ------------------- P1 (c) -------------------
\begin{paso}{1.c: Rapidez inicial $v_0$ del apoyo y posici\'on de encuentro $x^{*}$}
El apoyo parte en $t=2\,\mathrm{s}$. Con $\tau=t-2$ (su tiempo de movimiento) y $a=3$:
\[
v_{\text{apoyo}}(\tau)=v_0+3\tau,\qquad
x_{\text{apoyo}}(\tau)=v_0\,\tau+\tfrac12(3)\tau^2 .
\]

\textbf{Condici\'on 1 (rapidez en el encuentro).} El apoyo alcanza $v_f=26\,\mathrm{m/s}$:
\[
v_0+3\tau=26 \;\Rightarrow\; v_0=26-3\tau. \tag{i}
\]

\textbf{Condici\'on 2 (misma posici\'on en la etapa 2).} El encuentro ocurre en la
etapa 2 del cami\'on, donde $x_{\text{cam}}=96+24\,(t-8)=96+24(\tau-6)$:
\[
v_0\tau+1{,}5\,\tau^2 = 96+24(\tau-6). \tag{ii}
\]

Reemplazando (i) en (ii):
\[
(26-3\tau)\tau+1{,}5\tau^2 = 24\tau-48
\;\Longrightarrow\; 1{,}5\,\tau^2-2\tau-48=0
\;\Longleftrightarrow\; 3\tau^2-4\tau-96=0.
\]
\[
\tau=\frac{4+\sqrt{16+1152}}{6}=\frac{4+\sqrt{1168}}{6}\approx 6{,}36\ \mathrm{s}
\;\Rightarrow\; t^{*}=\tau+2\approx 8{,}36\ \mathrm{s}.
\]
\begin{justif}{Verificaci\'on de la etapa}
$t^{*}\approx 8{,}36\,\mathrm{s}\in[8,20]$: el encuentro efectivamente ocurre en la
\textbf{segunda etapa}, como pide el enunciado.
\end{justif}

Rapidez inicial y posici\'on de encuentro:
\[
v_0=26-3(6{,}36)\approx 6{,}91\ \mathrm{m/s},\qquad
x^{*}=96+24(8{,}36-8)\approx 104{,}70\ \mathrm{m}.
\]
\begin{resultado}\centering
$\displaystyle v_0 \approx 6{,}91\ \mathrm{m/s}\qquad\text{y}\qquad x^{*}\approx 104{,}70\ \mathrm{m}$
\end{resultado}
\end{paso}

\vspace{0.6cm}
\begin{center}
    \small\textbf{Usach Premium} --- ¡Nos vemos en la ayudant\'ia! \\
    \textit{"Por y para estudiantes."}
\end{center}

\end{document}
